\section{Reasoning Models, Exactly On-Policy (Run D)}\label{sec:rund}
The ProcessBench traces above are near on-policy, and their generators are instruction models.
Run D (preregistered, \texttt{docs/PREREGISTRATION\_RUN\_D.md}, with deviations recorded before
launch) asks the same questions of two long-reasoning models whose traces are exactly on-policy:
Qwen3-1.7B in thinking mode and DeepSeek-R1-Distill-Qwen-1.5B, on MATH-500 problems 21--500, with a
6,144-token budget per answer. Each model writes one trace per problem. Each failed trace is cut at
fixed token positions $t\in\{0,256,\dots,5120\}$, moved back to the last paragraph break so that a cut
depends only on the past, and $m=8$ continuations are sampled from each cut under the same total
budget (Fig.~\ref{fig:rund}).

\textbf{Failures are mostly budget exhaustion.} Accuracy is \rdQwenAcc\% (Qwen3) and \rdROneAcc\%
(R1-distill). 95\% and 92\% of their failures end because the reasoning hits the budget, not because
of a wrong final answer: these failures have no wrong step for a localiser to find.

\textbf{Value curves.} On 2$\times$T4 GPUs only R1-distill's stage~2 produced data in the session
limit: the first \rdROneN\ failures in dataset order. Proposition~1's prediction (D1) holds: no cut
is worth more than a fresh start, and $v(t)-v(0)$ falls from $0$ at $t=256$ to \rdDOneLast\
[\rdDOneLastLo, \rdDOneLastHi] at $t=5120$. \rdCap\% of these failures are capability-limited. At a
matched budget of generated tokens (D2), the preregistered estimator finds the best cut ahead of restart
by \rdDTwo\ [\rdDTwoLo, \rdDTwoHi] at 4$\times$, so the preregistered ``restart not beaten'' criterion
fails. The gain is fragile. Per attempt no cut beats restart, so it comes entirely from cheaper
attempts: a late cut leaves a short budget, and its attempts cost a fifth of a restart's. The estimator
extrapolates four observed samples to up to about 19 attempts, and its Jeffreys prior credits a cut
with no observed success with about 10\% per attempt. A plug-in estimate gives \rdDTwoPlug\
[\rdDTwoPlugLo, \rdDTwoPlugHi] (exploratory). On reasoning models, resampling late in a failed trace
may therefore pay through shorter attempts, as Proposition~2 allows, but our samples cannot settle
it. Budget-truncated traces also make part of the fall in $v(t)$ mechanical, since later cuts leave
less budget; only 4 wrong-answer failures have value curves (D4).

\begin{figure*}[t]\centering
\includegraphics[width=0.9\textwidth]{fig29_reasoning.pdf}
\caption{Run D, DeepSeek-R1-Distill-Qwen-1.5B, \rdROneN\ failed traces. (a) Value of a failed partial
trace cut at $t$ tokens, minus a fresh start (95\% bootstrap band). (b) Best cut (chosen on half of the
continuations, scored on the other) minus restart at a matched budget of generated tokens, under the
preregistered Jeffreys estimator and a plug-in sensitivity.}
\label{fig:rund}
\end{figure*}
